\chapter{Fallstudien}\label{scenarios}
Anhand der nachfolgenden Fallstudien und Szenarien werden die dieser Arbeit zugrundeliegenden Konzepte zur Errechnung eines Bauplanentwurfs aus einem als 3D Modell vorliegenden Gebäudeplans veranschaulicht und auf Anwendbarkeit überprüft.
Dabei werden die Szenarien zunehmend komplexer, um auch das Zusammenspiel verschiedener Teilkonzepte zur Lösung einzelner Probleme zu verifizieren.
Abschließend wird für einige Szenarien untersucht wie Regeln erstens definiert und zweitens auf deren Ergebnisse angewandt werden können.
Damit soll die ungeordnete Menge von konkreten Bausteinen eines Bauplanentwurfs durch Regeln in eine Reihenfolge gebracht werden, sodass daraus ein schrittweise umsetzbarer Bauplan entsteht.

\section{Von 3D Gebäudeplan zu Bauplanentwurf}
Diese erste Fallstudie enthält drei Szenarien anhand derer die Konzepte zur Errechnung eines Bauplanentwurfs ausgehend von 3D Gebäudeplänen erläutert und getestet werden.
Die Pläne sind das Ergebnis der Gebäudemodellierung innerhalb eines Konstruktionsplaners.
Die Modellierung soll mithilfe der in Kapitel~\ref{basics} näher behandelten Technologien geschehen, weshalb die Gebäudepläne in ihrer Struktur einem verbreiteten Industriestandard entsprechen.

\subsection{Szenario Turm}\label{scenarios:scenario1}
In diesem einleitenden Szenario werden mithilfe eines einfachen Turmes mit vier Wänden einige der Kernkonzepte geprüft.
Dabei werden sowohl die Modellierung des Turmes innerhalb des Konstruktionsplaners, als auch das anschließende Errechnen des Bauplanentwurfs thematisiert.
Insbesondere das Anwenden verschiedener Mauerwerksverbände soll die Flexibilität des erarbeiteten Vorgehens demonstrieren.
\subsubsection*{Beschreibung}
Der Turm besteht lediglich aus vier 20 Meter hohen Wänden, die einen einzigen Raum einschließen.
Er hat einen Grundriss von 10$\times$10 Metern.
Die Wände sollen daraufhin unter Anwendung folgender Mauerwerksverbände realisiert werden:
\begin{itemize}
  \item Einem Läuferverband mit einem Versatz von 50 \% der Bausteinlänge.
  \item Einem Kopf/Binderverband.
  \item Einem Kreuzverband.
\end{itemize}
Da die beiden letzten Verbände bei gleichbleibendem Modul die Wanddicke verdoppeln, muss das Modell etwas angepasst werden, um nach wie vor denselben Grundriss aufzuweisen.
Daher gibt es sowohl einen Bauplan mit 1 Meter dicken als auch einen mit 2 Meter dicken Wänden.
Als Basismodul wird ein Baustein mit den Maßen 2$\times$1$\times$0.5 Metern verwendet.
Das vereinfacht die Interpretation der generierten Lösungen.
Das dazugehörige Raster hat die Größe 1$\times$1$\times$0.5 Meter.

\subsubsection*{Problemstellungen}
Durch Lösen dieses Szenarios werden folgende Fragestellungen beantwortet:
\begin{itemize}
  \item Wie lässt sich die Modellierung des Gebäudes innerhalb des Konstruktionsplaners trotz Integration eines Industriestandards intuitiv gestalten?
  \item In welcher Art muss das Basismodul vorliegen und wie kann es notfalls angepasst werden? 
  Denn oftmals werden an Wandenden oder Ecken Bausteine benötigt, die eine geringere Länge aufweisen.
  Dies entspricht in Realität dem Zerschneiden von Bausteinen.
  \item Wie können einem Algorithmus die verschiedenen Mauerwerksverbände sinnvoll vorgegeben werden?
  \item Wie können Eckbereiche gefunden und der jeweilige Mauerwerksverband auch an diesen Stellen passend angebracht werden ohne das Überbindemaß (siehe Kapitel~\ref{basics:Mauerwerksverband}) zu verletzen?
  \item Wie sieht das Ergebnis des erarbeiteten Konzepts aus und welche Informationen werden in den resultierenden Bauplanentwurf integriert?
\end{itemize}

\subsection{Szenario LEGO Klemmbausteine}\label{scenarios:scenario2}
Nun folgt ein Szenario, das den Gegebenheiten eines realen Gebäudes bis auf dessen Skalierung eher entspricht.
Es soll ein LEGO-Gebäude konstruiert werden, welches einen Innenraum, Fenster und Türen enthält.
Darum wird das Modul stark verkleinert, sodass es den Maßen eines 4$\times$2 LEGO Steins entspricht (siehe Kapitel~\ref{basics:lego}).

\begin{figure}[ht]
  \centering
  \includegraphics[width=0.505\columnwidth]{fig/scenario2_storey_plan.png}
  \caption{Grundriss Szenario 2.}\label{fig:scenarios:Scenario2 Gebaeudeplan}
\end{figure}

\subsubsection*{Beschreibung}
Zu sehen ist der Plan eines einfachen Hauses mit einem Stockwerk.
Dieses besitzt eine Eingangstür, eine Terrassentür neben einem Fenster und eine Tür, die das Badezimmer vom Hauptraum trennt.
Türen und Fenster stellen eine Herausforderung für den Planungsalgorithmus dar, da der Verlauf einer ansonsten durchgängigen Wand dadurch unterbrochen wird und Lücken aufweist.
Es gibt breite Außen- und dünne Innenwände.
Dafür müssen zwei Wandtypen definiert werden, die jeweils unterschiedliche Wanddicken vorgeben.
Diese entsprechen in ihren Maßen dem Raster, welches das \textit{LEGO System} vorgibt.
Dickere Wände sollen zwei Noppen breit sein.
Daher wird für diesen Wandtyp ein Grundmodul mit den Maßen 31.8$\times$15.8$\times$9.6 Millimetern und einem Raster von 8$\times$8$\times$9.6 Millimetern gewählt.
Für die dünneren Innenwände soll eine Breite von einer Noppe verwendet werden.
Dies entspricht einem Grundmodul von 15.8$\times$7.8$\times$9.6 Millimetern mit einem gleichbleibenden Raster.
Mögliche Rotationen von Wänden, Fenstern und Türen sind auf 90\textdegree{} Schritte limitiert.
Das stellt im Fall von LEGO eine vertretbare Einschränkung dar, da es ohnehin dem intuitiven Umgang mit dessen Steinen und gleichzeitig dem Baustil der meisten einfachen Gebäuden entspricht.

\subsubsection*{Problemstellungen}
\begin{itemize}
  \item Wie werden Öffnungen im Konstruktionsplaner umgesetzt und im Anschluss weiter verarbeitet?
  \item Wände können verschiedene Dicken aufweisen. 
  Wie kann das während der Modellierung angegeben werden? 
  Wie kann das zu verwendende Modul und der Mauerwerksverband für eine Wand vorgegeben werden?
  \item Wie können Übergänge zwischen Wänden verschiedener Dicke oder unterschiedlichem Verband realisiert werden?
  %\item Wie können Eckbereiche zwischen zwei Wänden mit unterschiedlicher Höhe behandelt werden?
  \item In diesem Szenario gibt es T-Kreuzungen. 
  Wie können T-Kreuzungen von Eckbereichen unterschieden und dem Mauerwerksverband entsprechend umgesetzt werden? 
\end{itemize}

\subsection{Szenario 3}\label{scenarios:scenario3}
Um das Konzept nicht nur an eigens modellierten Gebäuden zu evaluieren, werden für dieses Szenario externe Modelle herangezogen.
\subsubsection*{Beschreibung}
Dabei kann sich zum Beispiel an den Beispielprojekten des \textit{Karlsruher Instituts für Technologie} (KIT) bedient werden~\cite{KITSAMPLEHOUSE:online}.
Zum Zeitpunkt der Erstellung dieser Arbeit existieren dort drei verschiedene Modelle.
Getestet wird das Konzept an den beiden Gebäuden aus Abbildung~\ref{fig:scenarios:Scenario3 KIT}.
\begin{figure}[htb]
  \begin{subfigure}[b]{0.5\columnwidth}
    \includegraphics[width=\columnwidth]{fig/KIT-Haus.png}
    \caption{Modell eines einfachen Hauses.}\label{fig:scenarios:scenario3 kit haus}
  \end{subfigure}
  \hfill
  \begin{subfigure}[b]{0.5\columnwidth}
    \includegraphics[width=\columnwidth]{fig/KIT-Buerogebaeude.png}
    \caption{Modell eines großen Bürogebäudes.}\label{fig:scenarios:scenario3 kit office}
  \end{subfigure}
  \caption{Beispielmodelle des Karlsruher Instituts für Technologie (KIT)~\cite{KITSAMPLEHOUSE:online}.}\label{fig:scenarios:Scenario3 KIT}
\end{figure}
\begin{comment}
\end{comment}
Darum muss es möglich sein, zu einem bereits abgeschlossenen Gebäudeplan weiter Informationen hinzuzufügen.
Dazu zählen wie in den vorherigen Szenarios etwa die Angabe verschiedener Wandtypen und die dafür gewünschten Grundmodule und Mauerwerksverbände.

\subsubsection*{Problemstellungen}
\begin{itemize}
  \item Wie lassen sich vorhandene Gebäudemodelle um die für diese Arbeit notwendigen Informationen erweitern?
  \item Oft sind Gebäudepläne sehr detailliert und gespickt mit Informationen aus allen Fachbereichen des Bauwesens. Wie können lediglich die relevanten Daten aus dem Modell gefiltert werden?
\end{itemize}

\section{Regelbasierte Bauplandeduktion aus einem Bauplanentwurf}\label{scenarios:scenario4}
In dieser Fallstudie wird untersucht wie das Integrieren von Regeln in das erarbeitete Konzept ermöglicht werden kann.
Dabei ist es wünschenswert Regeln erst im Nachgang an den vorangegangenen Schritt des Berechnens eines Bauplanentwurfs zu definieren und auf diesem anzuwenden.
Darum dürfen konkrete Regelbeschreibungen nicht Teil der Programmcodebasis dieser Arbeit sein, sondern müssen auf andere Weise in das Projekt eingefügt werden können.

\subsection*{Problemstellungen}
\begin{itemize}
  \item Wie können Regeln und Regelsets formal definiert werden?
  \item Wie können diese Regelsets zur Auswertung eines Bauplanentwurfs herangezogen werden?
\end{itemize}